\begin{resumo}[RESUMO]
\begin{SingleSpace}
\noindent A malha de atitude de um drone é um sistema de controle crítico: se a leitura do sensor, o cálculo ou a atuação atrasam, o veículo cai. Este trabalho reproduz esse problema em uma haste de um grau de liberdade, com motor e hélice em uma ponta e um sensor MPU6050 na outra, controlada por um STM32F103 com FreeRTOS. A malha de controle roda em uma tarefa acordada a cada 5~ms por um semáforo liberado pelo temporizador; comandos e telemetria ficam em tarefas de menor prioridade. Medido no microcontrolador, o atraso entre a interrupção e a atualização do PWM foi de no máximo 580~$\mu$s (11,6\% do período), sem nenhum ciclo perdido em 34\,870. Os principais problemas foram o aquecimento do motor, resolvido com um indutor em série; o ruído do acelerômetro com a hélice girando e a deriva do giroscópio, tratados com um filtro complementar que corrige o ângulo devagar durante o voo; e a oscilação do PID, tratada com sintonia por margem de fase, anti-\textit{windup} e desligamento do motor por ângulo. Na bancada, a haste pairou com erro quadrático médio de 0,28°, e um refino com cinco conjuntos de ganhos reduziu em 28\% o erro acumulado em degraus de \textit{setpoint}.
\par\vspace{2\baselineskip}
\noindent\textbf{Palavras-chave: Controle de posição. Sistemas de tempo real. Modelagem dinâmica. Simulação. Sistemas embarcados.}
\end{SingleSpace}
\end{resumo}
\begin{otherlanguage*}{english}
\begin{resumo}[ABSTRACT]
\begin{SingleSpace}
\noindent The attitude loop of a drone is a safety-critical control system: if sensing, computation or actuation is late, the vehicle falls. This work reproduces that problem on a one-degree-of-freedom rod, with a motor and propeller at one end and an MPU6050 sensor at the other, controlled by an STM32F103 running FreeRTOS. The control loop runs in a task woken every 5~ms by a semaphore released by the timer; commands and telemetry run in lower-priority tasks. Measured on the microcontroller, the delay from the timer interrupt to the PWM update was at most 580~$\mu$s (11.6\% of the period), with no lost cycle in 34,870. The main problems were motor heating, solved with a series inductor; accelerometer noise with the propeller running and gyroscope drift, handled by a complementary filter that corrects the angle slowly in flight; and PID oscillation, handled by phase-margin tuning, anti-windup and an angle cutoff for the motor. On the bench, the rod hovered with a 0.28° RMS error, and a refinement with five gain sets reduced the accumulated error in setpoint steps by 28\%.
\par\vspace{2\baselineskip}
\noindent\textbf{Keywords: Position control. Real-time systems. Dynamic modeling. Simulation. Embedded systems.}
\end{SingleSpace}
\end{resumo}
\end{otherlanguage*}
